\documentclass[11pt,a4paper]{article}
\usepackage{amsmath}
\usepackage{amssymb}
\usepackage{amsthm}
\usepackage{latexsym}
\usepackage{pstricks}
\usepackage{amsbsy}
\usepackage{amscd}
\usepackage{mathrsfs}
\usepackage{tipa}
\usepackage{txfonts}
\usepackage{amsfonts}
\usepackage{graphicx}
\usepackage{listings}
\usepackage{amsmath}
\usepackage{amssymb}
\usepackage[all]{xy}
\usepackage{enumitem}
\setenumerate[1]{itemsep=1pt,partopsep=0pt,parsep=\parskip,topsep=6pt}
\setitemize[1]{itemsep=1pt,partopsep=0pt,parsep=\parskip,topsep=6pt}
\setdescription{itemsep=1pt,partopsep=0pt,parsep=\parskip,topsep=6pt}
\input xypic
\addtolength{\textheight}{1cm}
\addtolength{\textwidth}{3.8cm}
\addtolength{\oddsidemargin}{-1,8cm}
\addtolength{\evensidemargin}{-1,8cm}
\newtheorem{thm}{Theorem}[section]
\newtheorem{cor}[thm]{Corollary}
\newtheorem{lem}[thm]{Lemma}
\newtheorem{cond}{Condition}
\newtheorem{que}[thm]{Question}
\newtheorem{exm}[thm]{Example}
\newtheorem{prop}[thm]{Proposition}
\newtheorem{defn}[thm]{Definition}
\newtheorem{rem}[thm]{Remark}
\newtheorem{nota}[thm]{Notation}
\parindent15pt

\begin{document}


\title{\Large \bf  Virtually Gorenstein Artin algebras are weakly Gorenstein
\footnotetext{
*Corresponding author.  E-mail address: sdwg001@163.com.\\
This research was supported by the Natural Science Foundation of Hunan Province (No. 2023JJ30561).
 }}



 \vskip 0.5cm
\author{{\normalsize  Weiqing Li }\\
{\small College of Mathematics and Statistics, Jishou University, Jishou 416000,  P. R. China}\date{}}
 \maketitle
 \vskip-1cm


\begin{abstract}  We prove that  all semi-Gorenstein-projective modules over a virtually Gorenstein Artin algebra  are Gorenstein projective;   this
answers affirmatively a question proposed by Ringel and  Zhang in \cite{rz1}.
It turns out that the Auslander-Gorenstein Conjecture, the (Strong) Nakayama Conjecture and Tachikawa's First Conjecture hold for virtually
Gorenstein Artin algebras. We also establish a new criterion for determining  the projective (injective) dimensions    of modules over virtually
Gorenstein Artin algebras. This allows us to show  the  validity  of the Auslander-Reiten  Conjecture for (infinitely generated) modules over  Artin algebras such that all the Gorenstein-projective modules are
projective.
\end{abstract}







\def\s{\stackrel}
\def\Longrightarrow{{\longrightarrow}}
\def\A{\mathcal{A}}
\def\B{\mathcal{B}}
\def\C{\mathcal{C}}
\def\D{\mathcal{D}}
\def\T{\mathcal{T}}
\def\R{\mathcal{R}}
\def\P{\mathcal{P}}
\def\S{\mathcal{S}}
\def\H{\mathcal{H}}
\def\U{\mathscr{U}}
\def\V{\mathscr{V}}
\def\M{\mathscr{M}}
\def\N{\mathcal{N}}
\def\W{\mathscr{W}}
\def\X{\mathscr{X}}
\def\Y{\mathscr{Y}}
\def\Z{\mathcal {Z}}
\def\I{\mathcal {I}}
\def\add{\mbox{add}}
\def\Aut{\mbox{Aut}}
\def\coker{\mbox{coker}}
\def\deg{\mbox{deg}}
\def\diag{\mbox{diag}}
\def\dim{\mbox{dim}}
\def\End{\mbox{End}}
\def\Ext{\mbox{Ext}}
\def\Hom{\mbox{Hom}}
\def\Gr{\mbox{Gr}}
\def\id{\mbox{id}}
\def\Im{\mbox{Im}}
\def\ind{\mbox{ind}}
\def\mod{\mbox{mod}}
\def\mul{\multiput}
\def\c{\circ}
\def \text{\mbox}

\hyphenation{ap-pro-xi-ma-tion}



\textbf{Key words:}    Virtually
Gorenstein Artin algebra;  weakly Gorenstein Artin algebra; Auslander-Gorenstein Conjecture;   (Strong) Nakayama Conjecture; Auslander-Reiten  Conjecture.

\medskip

\textbf{2020 Mathematics Subject Classification:}  16G10; 18G25.



\section{ Introduction}



Throughout this paper, $\Lambda$ is an Artin algebra and
all $\Lambda$-modules are unitary  right  $\Lambda$-modules. We view left $\Lambda$-modules as right  ones over
the opposite ring $\Lambda^{{\rm op}}$. We denote by (mod-$\Lambda$)  Mod-$\Lambda$ the class of all  (finitely generated) right    $\Lambda$-modules.
The symbols  id$_{\Lambda}$$(M)$ and pd$_{\Lambda}$$(M)$  stand  for the   injective and projective   dimension of
  a  right $\Lambda$-module $M$, respectively.

The class of Gorenstein projective  modules plays a central role in the theory of Gorenstein homological algebra.

\begin{defn} \label{defn1} \rm{  (\cite{eno9}).
 A  {\it  complete projective resolution} is an exact sequence of projective
 $\Lambda$-modules,  $$\cdots  \rightarrow P_{1} \rightarrow P_{0}\rightarrow
P_{-1}\rightarrow P_{-2}\rightarrow \cdots,  $$
 such that Hom$_{\Lambda}$($-$, $Q$) leaves the sequence exact whenever $Q$ is a projective   $\Lambda$-module; the module $M$ = im($P_{0}\rightarrow
P_{-1}$) is then  said to be {\it Gorenstein  projective}.

{\it Gorenstein  injective }  modules are defined dually.
 We use $\mathscr{GP}$ and $\mathscr{GI}$    to denote, respectively,  the classes of all Gorenstein  projective and  injective
   (right)  $\Lambda$-modules. }\end{defn}

\begin{defn} \label{defn81} \rm{   (\cite{rz1}).
 A right $\Lambda$-module
$M$   is called   {\it semi-Gorenstein-projective} if $M\in$ $^{\perp_{\infty}}{\Lambda}$, i.e.,    Ext$ _{\Lambda}^{i}(M, \Lambda)=0$  for all  $i \geq 1$.
 }\end{defn}


It is clear that $\mathscr{GP} \subseteq$  $^{\perp_{\infty}}{\Lambda}$; but we learn from   \cite{rz1}
that   there exist  semi-Gorenstein-projective modules which are not Gorenstein projective.

\begin{defn} \label{defn82} \rm{   (\cite{rz1}).
 An Artin algebra $\Lambda$  is said to be   {\it right  weakly Gorenstein} provided that $^{\perp_{\infty}}{\Lambda} $ $\cap$ mod-$\Lambda \subseteq$ $\mathscr{GP}$.
 If both  $\Lambda$ and $\Lambda^{{\rm op}}$ are right  weakly Gorenstein,   then $\Lambda$  is called {\it   weakly Gorenstein}. }\end{defn}

In 2020, Ringel and  Zhang \cite{rz1} proposed the following problem.

\begin{que} \label{que1} { \rm   (\cite[Question 9.2]{rz1})}.
Is  an Artin algebra $\Lambda$  right  weakly Gorenstein in case all the right Gorenstein  projective modules are   projective?\end{que}


We concentrate our investigation to  the       question  above over virtually Gorenstein  Artin algebras, which is a natural enlargement of the class of Gorenstein algebras.

\begin{defn}\label{defn21} { \rm   (\cite[Definition 8.1]{ber})}.   {\rm  An  Artin algebra $\Lambda$  is said to be}   virtually Gorenstein   {\rm provided that}   $\mathscr{GP}^\perp$ = $^\perp\mathscr{GI}$.
  \end{defn}

It is well-known that   an  Artin algebra $\Lambda$ is virtually Gorenstein  if and only if $\Lambda^{{\rm op}}$ is  virtually Gorenstein (see \cite[Theorem 8.7]{bel}).
In this note,  we prove the following theorem (see Theorem \ref{thm101}), which gives an affirmative answer to
  Question \ref{que1}.
 \begin{thm}\label{thm81} Let   $\Lambda$  be  a virtually Gorenstein  Artin algebra. Then  $^{\perp_{\infty}}{\Lambda} $ $=$  $\mathscr{GP}$.
 In particular,  $\Lambda$ is weakly Gorenstein.
 \end{thm}


Let
 $$ 0  \longrightarrow \Lambda_{\Lambda}
\longrightarrow
E_0 \longrightarrow E_1
\longrightarrow \cdots \longrightarrow E_i
\longrightarrow \cdots   $$ be a minimal  injective resolution of $\Lambda_{\Lambda}$.

In 1958,   Nakayama \cite{nak}   proposed a conjecture, which by results
of M$\ddot{\rm u}$eller \cite{mue}   is equivalent to the following:
\\
\\
{\bf Nakayama Conjecture (NC) }\ \ {\it  Let   $\Lambda$  be    an Artin algebra. If    $E_i$ is projective for all $i \geq 0$, then  $\Lambda$ is self-injective.}
\\
\par
 Let $\Lambda$  be    an Artin $R$-algebra over a commutative  Artin ring $R$ with radical $J$. We denote by $D$: Mod-$\Lambda$  $\rightarrow$ Mod-$\Lambda^{{\rm op}}$
 the usual duality which is given by  $D$ = Hom$_{R}$($-$, $\mathbb{E}$) where $\mathbb{E}$ is the injective envelope of $R/J$.
 Tachikawa \cite{tac} gave the following conjecture which is equivalent to (NC).
 \\
\\
{\bf Tachikawa Conjecture (TC) }\ \ {\it  Let   $\Lambda$  be    an Artin algebra. }
\begin{itemize}
\item[$({\rm T1})$]If ${\rm Ext_{\Lambda}}^{i}(D(\Lambda^{{\rm op}}),  \Lambda_{\Lambda})=0$ for all  $i \geq 1$, then $\Lambda$ is self-injective.
\item[$({\rm T2})$] Assume  $\Lambda$ is self-injective. If $M \in$  mod-$\Lambda$   and  ${\rm Ext_{\Lambda}}^{i}(M, M )=0$ for all  $i \geq 1$, then $M$ is projective.\end{itemize}

In 1975, Auslander and Reiten \cite{ar0}  proposed   another conjecture, of which (NC)
is a special case.
 \\
\\
{\bf Generalized Nakayama Conjecture (GNC)  }\ \ {\it  Let   $\Lambda$  be    an Artin algebra.
 If $S \in$  {\rm mod}-$\Lambda$  is simple and  ${\rm Ext_{\Lambda}}^{i}(S, \Lambda)=0$ for all  $i \geq 0$, then $S$ is zero}. {\rm(This is equivalent to the statement that,
 each indecomposable injective $\Lambda$-module occurs as the summand of some $E_{i}$.)}
\\
\par
It is shown by Auslander and Reiten \cite{ar0} that  (GNC)  is equivalent to
the following conjecture.
 \\
\\
{\bf Auslander-Reiten  Conjecture (ARC) }\ \ {\it  Let   $\Lambda$  be    an Artin algebra.
 If $M \in$  {\rm mod}-$\Lambda$   and  ${\rm Ext_{\Lambda}}^{i}(M, M\oplus\Lambda)=0$ for all  $i \geq 1$, then $M$ is projective.}
\\
\par
In 1990, Colby and Fuller \cite{cof}  proposed   the following conjecture, of which (GNC)
is a special case.
\\
\\
{\bf Strong Nakayama Conjecture (SNC)  }\ \ {\it  Let   $\Lambda$  be    an Artin algebra.
 If $M \in$  {\rm mod}-$\Lambda$    and  ${\rm Ext_{\Lambda}}^{i}(M, \Lambda)=0$ for all  $i \geq 0$, then $M$ is zero.}
\\
\par
Recall that an  Artin algebra is said to be  $Gorenstein$ provided that both {\rm id}$_{\Lambda}$($\Lambda$)
 and {\rm id}$_{\Lambda^{{\rm op}}}$($\Lambda^{{\rm op}}$) are finite. The Nakayama Conjecture is also   a special case  of the following conjecture posted by  Auslander and Reiten \cite{ar2} in 1994.
\\
\\
{\bf Auslander-Gorenstein Conjecture (AGC)  }\ \ {\it  Let   $\Lambda$  be    an Artin algebra. If    {\rm pd}$_{\Lambda}$$(E_{i}) \leq i$ for all $i \geq 0$, then  $\Lambda$ is  Gorenstein.}
\\
\par
The following   conjecture was posed by  Auslander and Reiten in \cite[p. 150]{ar1} (see also \cite[Conjecture (13)]{ars}).
\\
\\
{\bf Gorenstein Symmetry Conjecture (GSC) }\ \ {\it  Let   $\Lambda$  be    an Artin algebra   with {\rm id}$_{\Lambda}$$(\Lambda) < \infty$. Then {\rm id}$_{\Lambda^{{\rm op}}}$$(\Lambda^{{\rm op}}) < \infty$}.
\\
\par
For an Artin algebra $\Lambda$,   we learn from \cite[p. 121]{ar1}  and  \cite[Corollary 2]{hua} that {\rm id}$_{\Lambda}$$(\Lambda) \leq 1$ if and only if {\rm id}$_{\Lambda^{{\rm op}}}$$(\Lambda^{{\rm op}}) \leq 1$.
Later, Beligiannis \cite[Theorem 11.4]{bel}  proved that GSC holds for the class of  virtually Gorenstein Artin  algebras.

It is well known that all the conjectures stated above are formal consequences of the conjecture formulated below.
\\
\\
{\bf Finitistic Dimension Conjecture (FDC) }\ \ {\it  Let   $\Lambda$  be    an Artin algebra. Then} findim$\Lambda$ := sup\{pd$_{\Lambda}(M)$ \ $|$ \ $M \in$  mod-$\Lambda$ with pd$_{\Lambda}(M)    < \infty$\} $< \infty$.
\\
\par
We summarize the relation of each conjecture in the following implication for Artin algebras,

 $${\rm(GSC)} \Leftarrow{\rm (FDC)}\Rightarrow {\rm (GNC)}   \Lleftarrow\mspace{-6mu}\Rrightarrow {\rm (ARC)} \Leftrightarrow {\rm (SNC)} \Rightarrow {\rm (AGC)}\Rightarrow {\rm (NC)} \Lleftarrow\mspace{-6mu}\Rrightarrow {\rm (TC),  }  $$
 where the notation (p) $\Lleftarrow\mspace{-6mu}\Rrightarrow $ (q) means that all  Artin algebras satisfy (p) if and only if all  Artin algebras satisfy (q), while (p) $\Rightarrow$ (q)  means that   every Artin algebra that  satisfy (p) also satisfy (q).

Although numerous known cases have been established where one of these conjectures holds (see, for instance, [2-6, 9-14, 24-26,   28] and the references therein),
the full resolution of all these conjectures still remains open.


Perhaps surprisingly,  as an immediate consequence of Theorem \ref{thm81}, we get the  validity of the Strong Nakayama Conjecture for
(infinitely generated) modules over virtually Gorenstein Artin  algebras (see Corollary \ref{cor10}).

\begin{cor}\label{cor881}   Let  $\Lambda$ be a virtually Gorenstein Artin  algebra,  and    $M\in $  {\rm Mod}-$\Lambda$. If ${\rm Ext}_{\Lambda}^{n}(M, \Lambda)=0$  for
 all $n \geq 0$, then $M  $ is zero.\end{cor}


It follows  that,    (GNC),  (NC), (AGC) and (T1) both are true over      virtually Gorenstein Artin  algebras.
We also  show  the  validity  of the Auslander-Reiten  Conjecture for (infinitely generated) modules over  Artin algebras such that all the Gorenstein-projective modules are
projective  (see Corollary  \ref{cor201}).
Finally, we provide another proof of the Gorenstein Symmetry Conjecture for virtually Gorenstein Artin algebrs (see  Theorem \ref{thm102}).




We continue the introduction by recalling some basic definitions and notions
that we use in this paper. Given a $\Lambda$-module $M$ and a class $ \mathscr{C}$ of    $\Lambda$-modules, we denote

  \begin{itemize}
 \item  $^\perp {\mathscr{C}}=\{X: {\rm Ext_{\Lambda}^{1}}(X, C)=0$  for all $C\in
{\mathscr {C}}\}$,
 \item  $ {\mathscr {C}}^\perp=\{X: {\rm Ext_{\Lambda}^{1}}(C, X)=0$
for all $C\in {\mathscr {C}}\}$,
\item  $^{\perp_{\infty}} {\mathscr{C}}=\{X: {\rm Ext}_{\Lambda}^{i}(X, C)=0$  for all $C\in
{\mathscr {C}} $  and $i \geq 1$\},
 \item  $ {\mathscr {C}}^{\perp_{\infty}}=\{X: {\rm Ext}_{\Lambda}^{i}(C, X)=0 $
for all $C\in {\mathscr {C}}$  and $i \geq 1$\},
\item  add$(\mathscr {C}) =$   the class of all direct summands of
finite direct sums of modules in $ \mathscr{C}$,
\item${\bf Proj}(\Lambda)$
= the class of all projective right $\Lambda$-modules,
\item${\bf Inj}(\Lambda)$
= the class of all injective right $\Lambda$-modules.\end{itemize}


{\bf   (Pre)covers and (pre)envelopes.} \ \ Let $\mathscr{C}$    be a  class of modules in  Mod-$\Lambda$.
A  morphism $\phi: C\rightarrow M$ of Mod-$\Lambda$ with $C\in \mathscr{C}$ is called a  $\mathscr{C}$-$precover$ \cite{eno5} of $M$ if for any
morphism $f: C'\rightarrow M$ with $C'\in \mathscr{C}$,
there is a morphism $g : C'\rightarrow C$ such that $\phi g = f$. Moreover, if the only such $g$ are automorphisms of $C$ when
$C'= C$ and $f= \phi $, then the $\mathscr{C}$-precover $\phi$ is
called a $ \mathscr{C}$-$cover$. Dually, we have the
definitions of a  $\mathscr{C}$-preenvelope and a $\mathscr{C}$-envelope.
\\
\par
{\bf   Cotorsion  theory.}\ \ A pair $(\mathscr{X},\mathscr{Y})$  of subcategories of  Mod-$\Lambda$   is said to be  a  {\it cotorsion  theory} \cite{eno5} if ${\mathscr{X}}^\perp = {\mathscr{Y}}
 $ and $^\perp{\mathscr {Y}}$ = ${\mathscr {X}}$. A cotorsion theory (${\mathscr{X},\ \mathscr{Y}}$)
 is called $complete$   if every object in  Mod-$\Lambda$   has a   special  ${\mathscr{X}}$-precover
and  a   special  $\mathscr
{Y}$-preenvelope.  A
cotorsion theory (${\mathscr{X},\ \mathscr{Y}}$)
 is said to be  $ hereditary$   if whenever
 $0\rightarrow Y\rightarrow Y'\rightarrow
Y''\rightarrow 0$ is exact with $Y$, $Y'\in {\mathscr{Y}}$
then $Y''\in {\mathscr{Y}}$.

\medskip

\section{A characterization of    modules of  virtually finite  injective
  dimension}



In this section, we   give a characterization of  {\it modules of  virtually finite  injective
  dimension} \cite{bel} over  right Noetherian rings.


We need   the following    proposition  due to $\check{\rm S}$aroch and $\check{\rm S}$$\check{\rm t}$ov$\acute{\rm {\i}}$$\check{\rm c}$ek.
\begin{prop}\label{prop1} {\rm (\cite[Theorem 5.6]{ss})}. For any ring $\Lambda$, $(^\perp\mathscr{GI}$, $\mathscr{GI})$ is a hereditary complete
  cotorsion theory.   \end{prop}

We recall some definitions from \cite{eno5} and \cite{gbe}.

 \begin{defn} \label{defn2} \rm{Let $\mu$ be an ordinal and $\mathscr{A}$ = ($A_{\alpha} $ $|$ $\alpha \leq \mu$) be a sequence of modules.
  If $A_{0} $ = 0,  $A_{\alpha} \subseteq  A_{\alpha + 1}$ for all $\alpha < \mu$ and   $A_{\alpha} $ = $\bigcup_{\beta < \alpha} A_{\beta} $ for all
 limit ordinals $\alpha \leq \mu$, then the sequence $\mathscr{A}$  is called a {\it continuous {\rm (}well ordered{\rm )} chain of modules.}

   Let $M$ be a module and $\mathscr{X}$ be a class of modules. $M$ is  $\mathscr{X}$-{\it filtered}, provided that there is an ordinal $\mu$
   and  a continuous chain of modules, ($M_{\alpha} $ $|$ $\alpha \leq \mu$), consisting of submodules of $M$ such that $M$ = $M_{\mu} $,
    and each of the modules $M_{\alpha + 1}/M_{\alpha}  $  ($\alpha < \mu$) is isomorphic to an element of $\mathscr{X}$. Then chain ($M_{\alpha} $ $|$ $\alpha \leq \mu$)
   is called an $\mathscr{X}$-{\it filtration} of $M$. If $\mu$ if finite, then  $M$ is said to be {\it finitely} $\mathscr{X}$-{\it filtered}.}\end{defn}


\begin{nota}\label{nota11}  {\rm   Let   $\mathscr{X}$ be a class of modules. The class of all finitely  $\mathscr{X}$-filtered modules is denoted by
 filt($\mathscr{X}$)}.
\end{nota}

For convenience, we list the following two results.

\begin{lem}\label{lem2}  {\rm (\cite[Corollary 6.14]{gbe})}. Let   $\mathscr{X}$ be a set of  $\Lambda$-modules containing $\Lambda$.
Then the class  $^\perp(\mathscr{X}^\perp)$ consists of all direct summands of $\mathscr{X}$-filtered modules.
\end{lem}

\begin{lem}\label{lem21}  {\rm (\cite[Corollary 1.14(ii)]{poj})}. Let   $\mathscr{X}$ be a set of finitely presented  $\Lambda$-modules containing $\Lambda$.
Then the class   {\rm add(filt($\mathscr{X}$))} coincides with the class of all finitely presented modules in $^\perp(\mathscr{X}^\perp)$.
\end{lem}

The following homological lemma is a generalization of \cite[Theorem 7.3.4]{eno5}.

\begin{lem}\label{lem1}  {\rm (\cite[Lemma 3.8]{mao}.)} Let $M$ and $N$ be   $\Lambda$-modules,    $n$  a positive integer and $M$
the union of a continuous chain of submodules $(M_{\alpha} $ $|$ $\alpha \leq \kappa)$. If ${\rm Ext}_{\Lambda}^{n}(M_{0}, N)=0$ and ${\rm Ext}_{\Lambda}^{n}(M_{\alpha + 1}/M_{\alpha}, N)=0$
whenever  $\alpha + 1 < \kappa$\, then ${\rm Ext}_{\Lambda}^{n}(M, N)=0$ \end{lem}

To simplify the proof  of the next theorem, we give the following notation.
\begin{nota}\label{nota1}  {\rm  \  Let  $\Lambda$ be a right  Noetherian ring. Then there is a
family ($E_{j})_{j\in J}$ of indecomposable injective right  $\Lambda$-modules such that every injective right  $\Lambda$-module is the direct sum
of copies of the various $E_{j}$ (see the proof of \cite[Theorem 5.4.1]{eno5}). For each $E_{j}$, take a  projective resolution $$ \cdots \stackrel {f_{j2}} \longrightarrow
{F_{j1}}\stackrel {f_{j1}} \longrightarrow {F_{j0}}\stackrel {f_{j0}}
\longrightarrow E_{j}  \longrightarrow  0  \eqno (\spadesuit_{j})$$
 of  $E_{j}$, and let
 $$ 0  \longrightarrow \Lambda_{\Lambda}\stackrel {h_{0}}
\longrightarrow
E_0\stackrel {h_{1}} \longrightarrow E_1\stackrel {h_{2}}
\longrightarrow \cdots \eqno (\clubsuit)$$
  be a minimal  injective resolution of   $\Lambda_{\Lambda}$. Put $\mathscr{X}$ = \{im$(f_{ji})$, im$(h_{i})$ \ $|$\  $i\geq 0, j\in J $\}}.\end{nota}

Now  we give a characterization of modules of virtually finite injective dimension   over right
Noetherian rings.



\begin{thm}\label{thm1}  Let  $\Lambda$ be a right  Noetherian ring and keep the notation as above. Then
$(^\perp(\mathscr{X}^\perp)$, $\mathscr{X}^\perp)$ is a hereditary complete
  cotorsion theory. Moreover,  $^\perp\mathscr{GI}$ = $^\perp(\mathscr{X}^\perp)$, i.e.,
 $\mathscr{X}^\perp$ = $\mathscr{GI}$. \end{thm}

\begin{proof}
Since $\mathscr{X}$ is a set, $(^\perp(\mathscr{X}^\perp)$, $\mathscr{X}^\perp)$ is a   complete
  cotorsion theory by  \cite[Theorem 7.4.1]{eno5}. Next we prove that it is hereditary.

  Since  each im$(f_{ji})$ belongs to $^\perp(\mathscr{X}^\perp)$, for any $n \geq 1$, we deduce from the exact sequence $(\spadesuit_{j})$
that ${\rm Ext}_{\Lambda}^{n}($im$(f_{ji})$, $Y)=0$ for all $Y \in \mathscr{X}^\perp$. Therefore, im$(f_{ji}) \in$  $^{\perp_{\infty}}(\mathscr{X}^\perp)$.
On the other hand, every injective right  $\Lambda$-module is the direct sum
of copies of the various $E_{j}$. Hence   ${\bf Inj}(\Lambda) \subseteq $  $^{\perp_{\infty}}(\mathscr{X}^\perp)$.
But then we see from from the exact sequence $(\clubsuit)$
that  im$(h_{i}) \in$  $^{\perp_{\infty}}(\mathscr{X}^\perp)$.
   In addition, the class  $^\perp(\mathscr{X}^\perp)$ consists of all direct summands of $\mathscr{X}$-filtered modules by Lemma \ref{lem2}.
These observations  together with Lemma \ref{lem1} tell us that, ${\rm Ext}_{\Lambda}^{n}($$B$, $Y)=0$ for all $B \in {^\perp(\mathscr{X}^\perp)}$, $Y \in \mathscr{X}^\perp$,  and $n \geq 1$.
So the pair $(^\perp(\mathscr{X}^\perp)$, $\mathscr{X}^\perp)$   is hereditary by \cite[Lemma 5.24]{gbe}.

Now we prove that  $\mathscr{X}^\perp$ = $\mathscr{GI}$.
Since injective modules are contained in the thick class $^\perp\mathscr{GI}$ (cf. \cite[Lemma 5.24]{ss}), we see from the exact sequences $(\spadesuit_{j})$  and $(\clubsuit)$
that $\mathscr{X} \subseteq {^\perp\mathscr{GI}}$. Hence $\mathscr{X}^\perp \supseteq {(^\perp\mathscr{GI}})^\perp$ = $\mathscr{GI}$ (cf. Proposition \ref{prop1}).

It remains to show that $\mathscr{X}^\perp $   $\subseteq \mathscr{GI}$. For any $M \in \mathscr{X}^\perp $, there is an exact sequence
$0\rightarrow  K    \rightarrow \Lambda^{(\lambda)} \rightarrow  M\rightarrow 0   $  for some set $\lambda$. Then the sequence
$0\rightarrow  \Lambda^{(\lambda)}  \rightarrow E_{0}^{(\lambda)} \rightarrow {\rm(im(}h_{1}))^{(\lambda)}\rightarrow 0   $  is also exact.
 Construct a pushout of  $\Lambda^{(\lambda)}  \rightarrow E_{0}^{(\lambda)}$ along $\Lambda^{(\lambda)} \rightarrow  M$, we obtain a commutative diagram with exact
rows  and columns:


 $$\xymatrix{&&0\ar[d]&0\ar[d]&\\0\ar[r]& K\ar[r]
\ar@{=}[d]&\Lambda^{(\lambda)}\ar[d]\ar[r]
&M\ar[d]\ar[r]&0\\
0\ar[r]&  K\ar[r]
& E_{0}^{(\lambda)}\ar[d]\ar[r]
& Q\ar[d]\ar[r]&0\\
 & &{\rm im(}h_{1})^{(\lambda)}\ar[d]\ar@{=}[r]&{\rm im(}h_{1})^{(\lambda)}\ar[d]&&&\\
&&0&0}$$
Note that ${\rm Ext}_{\Lambda}^{1}($${\rm im(}h_{1})^{(\lambda)}$, $M)   \cong  $ $({\rm Ext}_{\Lambda}^{1}$${\rm(im(}h_{1})$, $M))^{\lambda} =  0$.
Thus  $Q \cong M \oplus {\rm im(}h_{1})^{(\lambda)}$, and we get   an epimorphism $E_{0}^{(\lambda)} \rightarrow  M$ with $E_{0}^{(\lambda)}  $ injective ($\Lambda$ is right  Noetherian).
Hence there exists a short exact sequence   $$0 \longrightarrow K_{0} \longrightarrow I_{0}  \stackrel {g_{0}}\longrightarrow M \longrightarrow 0$$
where  $g_{0}$ is an  injective cover.  Then $K_{0}\in {\bf Inj}(\Lambda)^\perp$ by \cite[Corollary 7.2.3]{eno5}. But   $M \in {\bf Inj}(\Lambda)^{\perp_{\infty}}$.
 Thus, one can see from the above short exact sequence that $K \in  {\bf Inj}(\Lambda)^{\perp_{\infty}}$. For any $i \geq 1$, we have an exact sequence
  $$0 = {\rm Ext}_{\Lambda}^{1}{\rm(im}(h_{i}), M)  \rightarrow  {\rm Ext}_{\Lambda}^{2}{\rm(im}(h_{i}), K_{0})  \rightarrow  {\rm Ext}_{\Lambda}^{2}{\rm(im}(h_{i}), I_{0})  = 0.$$
So  $  {\rm Ext}_{\Lambda}^{2}{\rm(im}(h_{i}), K_{0})   = 0$.  Now consider the short exact sequence
$0\rightarrow  {\rm im(}h_{i - 1})  \rightarrow E_{i - 1}\rightarrow {\rm im(}h_{i})\rightarrow 0   $
which  induces an exact sequence  $$0 = {\rm Ext}_{\Lambda}^{1}(E_{i - 1}, K_{0})  \rightarrow  {\rm Ext}_{\Lambda}^{1}{\rm(im}(h_{i-1}), K_{0})  \rightarrow  {\rm Ext}_{\Lambda}^{2}{\rm(im}(h_{i}), K_{0})  = 0.$$
Thus $  {\rm Ext}_{\Lambda}^{1}{\rm(im}(h_{i-1}), K_{0})   = 0$ for any $i \geq 1$. This implies  $K \in \mathscr{X}^{\perp}$.
Recursively, we construct an exact sequence of right $\Lambda$-modules
$$ \cdots \stackrel {g_{2}} \longrightarrow
I_1\stackrel {g_{1}} \longrightarrow I_0\stackrel {g_{0}}
\longrightarrow M  \longrightarrow  0  $$
with $I_i \in {\bf Inj}(\Lambda)$ and ker$(g_i) \in {\bf Inj}(\Lambda)^{\perp_{\infty}}$ ($i \geq 0$).
Therefore, $M $   $\in \mathscr{GI}$ by \cite[Lemma 17(2)]{gii}, and   $\mathscr{X}^\perp $   $\subseteq \mathscr{GI}$, as required. \end{proof}





\medskip


\section{ Virtually Gorenstein Artin algebras are weakly Gorenstein}

We begin with the following   remark.



\begin{rem}\label{rem21}  {\rm Let $\mathscr{X}$ be the set as in  Notation \ref{nota1}.  If    $\Lambda$ is an Artin algebra,
then we can always assume that $\mathscr{X} \subseteq $ mod-$\Lambda$.}\end{rem}

In view of the above remark, we get the following lemma.
\begin{lem}\label{lem22}   Let    $\Lambda$ be a virtually
Gorenstein Artin algebra and $\mathscr{X}$ be the set as in {\rm Notation \ref{nota1}}. Then, for any  $M\in $  {\rm mod}-$\Lambda$, there are two short exact sequences in  {\rm mod}-$\Lambda$
 $$0 \rightarrow M  \rightarrow Q  \rightarrow B \rightarrow 0 \ \ \ \ \ \ and  \ \ \ \ \ \ 0 \rightarrow M  \rightarrow G  \rightarrow C \rightarrow 0$$
 with  $Q  $ and  $C  $ belong to {\rm add(filt($\mathscr{X}$))},  $B\in $  $\mathscr{GP}$ and $G\in $  $\mathscr{GI}$. \end{lem}

 \begin{proof} Since $\Lambda$ is  virtually
Gorenstein, we have a  hereditary  cotorsion triple ($\mathscr{GP}$,  $^\perp\mathscr{GI}$, $\mathscr{GI}$).
 Hence, from \cite[Theorem 8.2(vi)]{bel}, we get two short exact sequences in  {\rm mod}-$\Lambda$
 $$0 \rightarrow M  \rightarrow Q  \rightarrow B \rightarrow 0 \ \ \ \ \ \ and  \ \ \ \ \ \ 0 \rightarrow M  \rightarrow G  \rightarrow C \rightarrow 0$$
 with  $Q  $ and  $C  $ belong to $^\perp\mathscr{GI}$,  $B\in $  $\mathscr{GP}$ and $G\in $  $\mathscr{GI}$.
Note that  $^\perp\mathscr{GI}$ = $^\perp(\mathscr{X}^\perp)$ by Theorem \ref{thm1}. It follows from   Lemma  \ref{lem21} that  $Q  $ and  $C  $ are contained in  {\rm add(filt($\mathscr{X}$))},
finishing the proof.\end{proof}

To determine the  (pro)jective dimensions for   modules over a  virtually
Gorenstein Artin algebra,  we need three auxiliary lemmas.



\begin{lem}\label{lem81}   Let    $\Lambda$ be a virtually
Gorenstein Artin algebra and   $M\in $  {\rm Mod}-$\Lambda$. If   $M \in $ ${D(\Lambda^{{\rm op}})}^{\perp_{\infty}}\ \cap $ $\mathscr{GP}^\perp$,
  then $M $ is injective.\end{lem}

 \begin{proof}
Since  $\Lambda$ is an Artin algebra, there is a
family ($E_{j})_{j\in J}$ of finitely generated injective right  $\Lambda$-modules such that every injective right  $\Lambda$-module is the direct sum
of copies of the various $E_{j}$.
    Note that  $D(\Lambda^{{\rm op}})  $ is
an injective cogenerator and  $M \in $ ${D(\Lambda^{{\rm op}})}^{\perp_{\infty}} $.  So $M \in $  ${E_{j}}^{\perp_{\infty}}$  for all  $j \in J$. It follows that     $M \in $  {\bf Inj}$(\Lambda)^{\perp_{\infty}}$.


 Let $\mathscr{S}$ be an irredundant finite set   of representatives of the simple modules in Mod-$\Lambda$.  For any    $S\in $  $\mathscr{S}$,  there is a short exact sequence
$0\rightarrow  S    \rightarrow Q \rightarrow  B\rightarrow 0   $ with  $Q\in $    add(filt($\mathscr{X}$))   and   $B\in $  $\mathscr{GP}$ by  Lemma  \ref{lem22}.
Then $Q$ is  isomorphic to a  direct summand of some module $Y$  in
 filt($\mathscr{X}$). According to the construction of $\mathscr{X}$ and the fact that $M \in $  {\bf Inj}$(\Lambda)^{\perp_{\infty}}$,  using dimension shfting, one can easily see that
there exists some positive integer $i_{S}$ such that   {\rm Ext}$ _{\Lambda}^{t}(Y, M)=0$  for all  $t \geq i_{S}$. Hence   {\rm Ext}$ _{\Lambda}^{t}(Q, M)=0$  for all  $t \geq i_{S}$.
But $M \in $  $\mathscr{GP}^\perp$ = $ {\mathscr {GP}}^{\perp_{\infty}}$. We then infer from the above short exact sequence that  {\rm Ext}$ _{\Lambda}^{t}(S, M)=0$  for all  $t \geq i_{S}$.

Since the set   $\mathscr{S}$ finite,   $i$ = sup\{$i_{S}$ \ $|$\   $S\in $  $\mathscr{S}$ \} $  < \infty$. In addition,  any $V \in$ {\rm mod}-$\Lambda$
 is  finitely  $\mathscr{S}$-filtered by \cite[Proposition
11.1]{and}. Thus, using dimension shfting, one can easily get that  {\rm Ext}$ _{\Lambda}^{\geq i}(V, \Lambda)=0$  for any
$V \in$ {\rm mod}-$\Lambda$. Therefore, {\rm id}$_{\Lambda}$($M$) $\leq i <\infty$.
But   $M \in $  {\bf Inj}$(\Lambda)^{\perp_{\infty}}$. Therefore, using dimension shifting, one can easily check that  $M $ is injective.  This   finishes  the proof.
 \end{proof}

\begin{lem}\label{lem82}   Let    $\Lambda$ be a virtually
Gorenstein Artin algebra and   $M\in $  {\rm Mod}-$\Lambda$. If   $M \in $ ${D(\Lambda^{{\rm op}})}^{\perp_{\infty}} $,
  then $M $   $\in $  $\mathscr{GI}$.\end{lem}

\begin{proof} By Proposition \ref{prop1}, there is a short exact sequence
$0\rightarrow  K_{0}    \stackrel {g_{0}}\longrightarrow I_{0} \rightarrow  M\rightarrow 0   $ with  $I_{0} \in $    $^{\perp_{\infty}}{\mathscr{GI}}$   and   $ K_{0} \in $  $\mathscr{GI} \subseteq$ ${\bf Inj}(\Lambda)^{\perp_{\infty}}$.
Then $I_{0} \in $ ${D(\Lambda^{{\rm op}})}^{\perp_{\infty}} $ since $M \in $ ${D(\Lambda^{{\rm op}})}^{\perp_{\infty}} $.
But $^{\perp}{\mathscr{GI}}$ = $\mathscr{GP}^\perp$ as $\Lambda$ is virtually
Gorenstein. Hence  $I_{0} \in $ ${D(\Lambda^{{\rm op}})}^{\perp_{\infty}}\ \cap $ $\mathscr{GP}^\perp$.
It follows from  Lemma  \ref{lem81} that $I_{0}   $ is injective.

Recursively, we construct an exact sequence of right $\Lambda$-modules
$$ \cdots \stackrel {g_{2}} \longrightarrow
I_1\stackrel {g_{1}} \longrightarrow I_0\stackrel {g_{0}}
\longrightarrow M  \longrightarrow  0  $$
with $I_i \in {\bf Inj}(\Lambda)$ and ker$(g_i) \in {\bf Inj}(\Lambda)^{\perp_{\infty}}$ ($i \geq 0$).
Therefore, $M $   $\in \mathscr{GI}$ by \cite[Lemma 17(2)]{gii}, as required. \end{proof}

\begin{lem}\label{lem83}   Let    $\Lambda$ be a virtually
Gorenstein Artin algebra and   $M\in $  {\rm Mod}-$\Lambda$. If   $M \in $ ${\mathscr{GI}}^{\perp_{\infty}} $,
  then $M $   is injective.\end{lem}

\begin{proof} Since $M \in $ ${\mathscr{GI}}^{\perp_{\infty}} \subseteq$ ${D(\Lambda^{{\rm op}})}^{\perp_{\infty}} $,
  $M $   $\in \mathscr{GI}$ by Lemma \ref{lem82}.
So  $E(M)/M  $ $\in \mathscr{GI}$, where  $E(M)$ is the injective envelope of $M $.
 It follows that $M $ is a direct summand of $E(M)$ since  Ext$_{\Lambda}^{1}(E(M)/M, M) =  0$,  as desired. \end{proof}

We now establish a new criterion for determining the (Gorenstein) injective dimensions of   modules over virtually
Gorenstein Artin algebras.

 \begin{thm}\label{thm91}   Let    $\Lambda$ be a virtually
Gorenstein Artin algebra.  For any   $M\in $  {\rm Mod}-$\Lambda$, the following  two equalities hold.
 \begin{itemize}
\item[$(1)$]  {\rm  id$_{\Lambda}$$(M)$ =  sup\{$n \in \mathbb{N}$\ $|$\   Ext$_{\Lambda}^{n}(G, M) \neq  0$ for all $G $   $\in $  $\mathscr{GI}$\}.}
\item[$(2)$] {\rm  Gid$_{\Lambda}$$(M)$ =  sup\{$n \in \mathbb{N}$\ $|$\   Ext$_{\Lambda}^{n}(D(\Lambda^{{\rm op}}), M) \neq  0$\}.}
\end{itemize}    \end{thm}

\begin{proof} By virtue of Lemma \ref{lem83}, it is straightforward to adapt the argument presented in  \cite[Lemma 5.1]{bel1} to derive the equality  (1); the proof of (2) is similar
(but using Lemma \ref{lem82}).
\end{proof}

Next we prove the dual of Theorem \ref{thm91}.
 \begin{thm}\label{thm92}   Let    $\Lambda$ be a virtually
Gorenstein Artin algebra.  For any   $M\in $  {\rm Mod}-$\Lambda$, the following  two equalities hold.
 \begin{itemize}
\item[$(1)$]  {\rm  pd$_{\Lambda}$$(M)$ =  sup\{$n \in \mathbb{N}$\ $|$\   Ext$_{\Lambda}^{n}(M, Q) \neq  0$ for all $Q $   $\in $  $\mathscr{GP}$\}.}
\item[$(2)$] {\rm  Gpd$_{\Lambda}$$(M)$ =  sup\{$n \in \mathbb{N}$\ $|$\   Ext$_{\Lambda}^{n}(M, \Lambda) \neq  0$\}.}
\end{itemize}    \end{thm}



\begin{proof} We only prove  (1);  the proof of (2) is similar.

Write  $t$ := sup\{$n \in \mathbb{N}$\ $|$\   Ext$_{\Lambda}^{n}(M, Q) \neq  0$ for all $Q $   $\in $  $\mathscr{GP}$\}.
It is   clear that $t \leq$   pd$_{\Lambda}$$(M)$. Next we prove that    pd$_{\Lambda}$$(M)$ $\leq t $. We may
assume $t < \infty$. Let  $G $ be   an arbitrary Gorenstein injective right $\Lambda^{{\rm op}}$-module.
Then $D(G)  \in \mathscr {GP}$  by \cite[Lemma 8.6(1)]{bel}. Thus
$${\rm Ext}_{\Lambda^{{\rm op}}}^{t+l}(G, D(M)) \cong  D{\rm Tor}_{t+l}^{\Lambda}(M, G)  \cong      {\rm Ext}_{\Lambda}^{t+l}(M, D(G)) =  0$$ for all $l \geq 1 $
by  \cite[ Lemma 2.16(b)]{gbe}, and hence     id$_{\Lambda^{{\rm op}}}$$(D(M)) \leq t$   per Theorem \ref{thm91}. So  pd$_{\Lambda}$$(DD(M)) \leq t$.
In addition, $M  $ is a pure submodule of $DD(M)$ by \cite[Proposition 5.3.9]{eno5}. Therefore,  pd$_{\Lambda}$$(M)$ $\leq t $. This forces that  pd$_{\Lambda}$$(M)$ $= t $.
 \end{proof}

We see from the following result that the Strong Nakayama Conjecture is valid for (infinitely generated) modules over virtually Gorenstein Artin algebras.
\begin{cor}\label{cor10}   Let    $\Lambda$ be a virtually
Gorenstein Artin algebra.  For any   $M\in $  {\rm Mod}-$\Lambda$, if   ${\rm Ext_{\Lambda}}^{i}(M, \Lambda)=0$ for all  $i \geq 0$, then $M$ is zero.    \end{cor}
\begin{proof} This  is a direct consequence of Theorem \ref{thm92}.\end{proof}


\begin{rem}\label{rem29}  {\rm
From the previous corollary,  we get    the validity
of the (Generalized) Nakayama Conjecture over all  virtually
Gorenstein Artin algebras, refining a main result in \cite[Theorem 1.2]{ch1}.}\end{rem}

We are in a position to show that every virtually Gorenstein  Artin algebra is weakly Gorenstein.

 \begin{thm}\label{thm101} Let   $\Lambda$  be  a virtually Gorenstein  Artin algebra. Then  $^{\perp_{\infty}}{\Lambda} $ $=$  $\mathscr{GP}$.
 In particular,  $\Lambda$ is weakly Gorenstein.
 \end{thm}

 \begin{proof} Since $\Lambda$ is a virtually Gorenstein  Artin algebra, we obtain from Theorem \ref{thm92} that $^{\perp_{\infty}}{\Lambda} $ $=$  $\mathscr{GP}$.
 So $\Lambda$ is right weakly Gorenstein. But $\Lambda^{{\rm op}}$ is also   virtually Gorenstein (cf. \cite[Theorem 8.7]{bel}).
 Thus $\Lambda^{{\rm op}}$ is also right weakly Gorenstein. Therefore,  $\Lambda$ is weakly Gorenstein.\end{proof}

 We now give an affirmative answer to the question proposed by   Ringel and  Zhang in \cite[Question 9.2]{rz1}.
 \begin{cor}\label{cor101} Let   $\Lambda$  an  Artin algebra with ${\bf Proj}(\Lambda)$ $=$  $\mathscr{GP}$. Then   $\Lambda$ is weakly Gorenstein.
 \end{cor}

  \begin{proof} Note that  ${\bf Proj}(\Lambda)$ $=$  $\mathscr{GP}$. So $\mathscr{GP}^\perp$  = Mod-$\Lambda$, and hence $\mathscr{GP}^\perp $ $\cap$ mod-$\Lambda$  = mod-$\Lambda$
  is {\it contravariantly  finite} (in the sense of \cite{ar1}) in mod-$\Lambda$.
    We then infer from \cite[Theorem 8.2(viii)]{bel} that
  $\Lambda$ is  virtually Gorenstein.    Therefore,  $\Lambda$ is weakly Gorenstein by Theorem \ref{thm101}.\end{proof}

Surprisingly, we see from the next result that the  Auslander-Gorenstein Conjecture (AGC) is true  for (infinitely generated) modules over an Artin algebra  $\Lambda$  with ${\bf Proj}(\Lambda)$ $=$  $\mathscr{GP}$.
 \begin{cor}\label{cor201} Let   $\Lambda$  an  Artin algebra with ${\bf Proj}(\Lambda)$ $=$  $\mathscr{GP}$. If $M \in$  {\rm Mod}-$\Lambda$   and  ${\rm Ext_{\Lambda}}^{i}(M, \Lambda)=0$ for all  $i \geq 1$, then $M$ is projective.
 \end{cor}

  \begin{proof} Since  ${\bf Proj}(\Lambda)$ $=$  $\mathscr{GP}$,
  $\Lambda$ is  virtually Gorenstein.   Now  one can apply Theorem \ref{thm92} (2) to get that $M$ $\in \mathscr{GP}$ = $ {\bf Proj}(\Lambda)$.\end{proof}


 The following theorem provides
another proof of the Gorenstein Symmetry Conjecture for virtually Gorenstein Artin algebrs (see \cite[Theorem  11.4]{bel}).

\begin{thm}\label{thm102}  Let  $\Lambda$ be a virtually Gorenstein Artin algebra and  $ E_{i} $  $(i \geq 0)$ be the module as in  {\rm Notation \ref{nota1}}. Then the following quantities are identical.
\begin{itemize}
\item[$(1)$]  {\rm id$_{\Lambda}$($\Lambda$)}.
 \item[$(2)$] {\rm id$_{\Lambda^{{\rm op}}}$($\Lambda^{{\rm op}}$)}.
 \item[$(3)$]  {\rm sup\{$n \in \mathbb{N}$\ $|$\   Ext$_{\Lambda}^{n}(D(\Lambda^{{\rm op}}), \Lambda) \neq  0$\}}.
 \item[$(4)$]  {\rm sup\{$n \in \mathbb{N}$\ $|$\   Ext$_{\Lambda^{{\rm op}}}^{n}(D(\Lambda), \Lambda^{{\rm op}}) \neq  0$\}}.
  \item[$(5)$]  {\rm sup\{pd$_{\Lambda}$$(E_{i})$\ $|$\ $i \in \mathbb{N}$\}}.\end{itemize}\end{thm}

 \begin{proof} We get from Theorem \ref{thm91} that {\rm Gid$_{\Lambda}$($\Lambda$)} =   (3).
 But {\rm Gid$_{\Lambda}$($\Lambda$)} = {\rm id$_{\Lambda}$($\Lambda$)} per \cite[Theorem  2.1]{hol1}. So (1) $= (3)$.
Similarly, we have (2) $= (4)$.

Dually, we obtain (3) =   pd$_{\Lambda^{{\rm op}}}$$(D(\Lambda^{{\rm op}})$) by using Theorem \ref{thm92} and  \cite[Theorem  2.2]{hol1}.
 Note that $D(\Lambda^{{\rm op}})$ is an injective cogenerator. These observations force that (5) $\leq (3)$.


 Next we prove the equality (3) $= (4)$. Note that $DD(\Lambda) \cong \Lambda$  and $DD(\Lambda^{{\rm op}}) \cong \Lambda^{{\rm op}}$. Hence
${\rm Ext}_{\Lambda}^{n}(D(\Lambda^{{\rm op}}), \Lambda) \cong {\rm Ext}_{\Lambda}^{n}(D(\Lambda^{{\rm op}}), DD(\Lambda)) \cong
 D{\rm Tor}_{n}^{\Lambda}(D(\Lambda^{{\rm op}}), D(\Lambda))
 \cong      {\rm Ext}_{\Lambda^{{\rm op}}}^{n}(D(\Lambda), DD(\Lambda^{{\rm op}})) \cong      {\rm Ext}_{\Lambda^{{\rm op}}}^{n}(D(\Lambda), \Lambda^{{\rm op}})$
by  \cite[ Lemma 2.16(b)]{gbe}. Consequently, (3) $= (4)$.

 It remains to argue that (3) $\leq (5)$.  We may assume {\rm sup\{pd$_{\Lambda}$$(E_{i})$\ $|$\ $i \in \mathbb{N}$\}} = $t < \infty$.
  Let    $\mathscr{S}$  be an irredundant finite set of representatives of the simple modules in Mod-$\Lambda$.
Write $G = \bigoplus_{S\in \mathscr{S} } E(S)$, where $E(S)$ the injective envelope of    $S$.
 For any  $S\in$ $\mathscr{S}$,  we get some  $n \geq 0$ such that  ${\rm Ext}_{\Lambda}^{n}(S, \Lambda)\neq 0$
by Corollary \ref{cor10}.
 So $E(S)$ is a direct summand of the module $E_{n}$ in Notation \ref{nota1} (cf. \cite[p. 70]{ar0}), and hence {\rm pd}$_{\Lambda}$($E(S)) \leq  t$.
 Thus   pd$_{\Lambda}$$(G) \leq t$. But  $G $ is an injective cogenerator. Therefore,
  {\rm pd}$_{\Lambda}$($D(\Lambda^{{\rm op}}))$  $\leq t$, which implies   that (3) $\leq (5)$. The proof is finished.
 \end{proof}


\begin{rem}\label{rem30}  {\rm
The validity of both the Auslander-Gorenstein Conjecture and Tachikawa's First Conjecture in the setting of virtually Gorenstein Artin algebras
can be directly deduced as a special instance of Theorem \ref{thm102}.}\end{rem}


 In view of the above discussions, it is natural to close the present paper by putting forward the following conjecture,
  which can be regarded as a   generalized extension covering the Auslander-Gorenstein Conjecture as well as   Tachikawa's First Conjecture.
  \\
\\
{\bf Conjecture 3.15.} {\it Let  $\Lambda$ be an Artin algebra $\Lambda$ and  $ E_{i} $  $(i \geq 0)$ be the module as in  {\rm Notation \ref{nota1}}.
 Then the following quantities are identical.
\begin{itemize}
\item[$(1)$]  {\rm id$_{\Lambda}$($\Lambda$)}.
 \item[$(2)$]  {\rm sup\{$n \in \mathbb{N}$\ $|$\   Ext$_{\Lambda}^{n}(D(\Lambda^{{\rm op}}), \Lambda) \neq  0$\}}.
 \item[$(3)$]  {\rm sup\{pd$_{\Lambda}$$(E_{i})$\ $|$\ $i \in \mathbb{N}$\}}.\end{itemize}}


\section*{Acknowledgments}


I would like to thank Professor Hongxing Chen for many useful discussions on this paper.

\section*{Disclosure statement}
The author reports there are no competing interests to declare. This note has no associated data.

   \vskip3mm

\begin{thebibliography}{99}

\bibitem{and} F. W. Anderson,  K. R. Fuller,  {\it Rings and
Categories  of  Modules}, 2nd. Springer-Verlag, New York, 1992.

\bibitem{ar0}  M. Auslander, I. Reiten, On a generalized version of Nakayama conjecture, {\it Proc. Amer. Math. Soc.}  52 (1975), 69-74.

\bibitem{ar1}  M. Auslander, I. Reiten, Applications of contravariantly finite subcategories, {\it Adv. Math.} 86 (1991), 111-152.

\bibitem{ar2}  M. Auslander, I. Reiten,  $k$-Gorenstein algebras and syzygy modules, {\it J. Pure Appl.
Algebra} 92 (1994), 1-27.

\bibitem{ars} M. Auslander, I. Reiten, S. O. Smal${\o}$, {\it Representation Theory of Artin Algebras.} Cambridge Studies in Advanced Mathematics, vol. 36. Cambridge Univ.  Press, Cambridge, 1995.

\bibitem{bel} A. Beligiannis, Cohen-Macaulay modules, (co)torsion pairs and virtually Gorenstein algebras, {\it J. Algebra }  288 (2005), 137-211.

\bibitem{bel1} A. Beligiannis, On algebras of finite Cohen-Macaulay type,  {\it Adv. Math. }  226 (2011), 1973-2019.

\bibitem{ber} A. Beligiannis, I. Reiten,  Homological and homotopical aspects of torsion theories, {\it Mem. Amer. Math. Soc}. 188(883) (2007), 1-207.

\bibitem{cum} C. Cummings, Left-right symmetry of finite finitistic dimension, {\it Bull. Lond. Math. Soc.} 56(2) (2024),  624-633.


\bibitem{cfx} H. X. Chen, M. Fang, C. C. Xi, Tachikawa¡¯s second conjecture, derived recollements, and gendo-symmetric algebras, {\it Compos. Math.} 160(11),  (2024), 2704-2734.

\bibitem{cof} R. R. Colby, K. R. Fuller,  A note on the Nakayama Conjecture, {\it Tsukuba J. Math.}
14 (1990), 343-352.

\bibitem{chh} L. W. Christensen, H. Holm,  Algebras that satisfy Auslander's condition on vanishing
of cohomology, {\it Math. Z.} 265 (2010),  21-40.


\bibitem{chx} H. X. Chen, C. C. Xi,  Homological theory of self-orthogonal modules,  {\it Trans. Amer. Math. Soc.} 378(10) (2025),  7287-7335.

\bibitem{ch1} H. X. Chen,   C. C. Xi,  Virtually Gorenstein algebras of infinite dominant dimension, {\it J. Pure Appl.
Algebra} 230(4) (2026), 108224.

\bibitem{eno9}  E. E. Enochs, O. M. G. Jend, Gorenstein injective and projective modules,  {\it Math. Z}. 220(4)  (1995), 611-633.

\bibitem{eno5} E. E. Enochs, O. M. G. Jend, {\it Relative  Homological  Algebra},  Berlin,  Walter de Gruyter,
2000.


\bibitem{gii} J. Gillespie, A. Iacob, Duality pairs, generalized Gorenstein modules, and Ding injective
envelopes, {\it Comptes Rendus Math$\acute{{\rm e}}$matique}  360 (2022), 381-398.

\bibitem{gbe} R. G$\ddot{\rm o}$bel, J. Trlifaj, {\it Approximations and Endomorphism Algebras of Modules}, de Gruyter Exp. Math. vol.
41,  2 revised and extended edn. de Gruyter, Berlin, (2012).

\bibitem{hol1} H. Holm,  Rings with finite gorenstein injective dimension, {\it Proc. Amer. Math. Soc.} 132(5) (2003), 1279-1283.


\bibitem{hua} Z. Y. Huang, Selforthogonal modules with finite injective dimension II, {\it J. Algebra} 264 (2003),   262-268.

\bibitem{mao}  L. X. Mao,  N. Q. Ding,  Relative cotorsion modules and relative flat modules, {\it Comm. Algebra} 34 (2006), 2303-2317.

\bibitem{mue} B. J.  M$\ddot{\rm u}$eller, The classification of algebras by dominant dimension, {\it Canad. J. Math.} 20 (1968), 398-409.

 \bibitem{poj} P. Moradifar, J. $\check{\rm S}$aroch, J. $\check{\rm S}$$\check{\rm t}$ov$\acute{\rm {\i}}$$\check{\rm c}$ek,   Finitistic dimension conjectures via Gorenstein
projective dimension, {\it J. Algebra} 591 (2022), 15-35.

\bibitem{nak} T. Nakayama, On algebras with complete homology, {\it Abh. Math. Sem. Univ. Hamburg} 22  (1958), 300-307.

\bibitem{rz1} C. M. Ringel, P. Zhang, Gorenstein-projective and semi-Gorenstein-projective modules, {\it Algebra $\&$   Number Theory}
14(1) (2020), 1-36.

\bibitem{riz}  C. M. Ringel, P. Zhang, Gorenstein-projective modules over short local algebras. {\it J. Lond. Math. Soc.} 106(2) (2022), 528-589.

 \bibitem{ss} J. $\check{\rm S}$aroch, J. $\check{\rm S}$$\check{\rm t}$ov$\acute{\rm {\i}}$$\check{\rm c}$ek,   Singular compactness and definability for $\sum$-cotorsion and Gorenstein modules,

{\it Sel. Math. New Ser.} 26:23 (2020).

\bibitem{tac} H. Tachikawa, {\it Quasi-Frobenius rings and generalizations:  QF-3 and QF-1 rings}, Notes by Claus Michael Ringel. Lecture Notes in Mathematics, 351, Berlin-Heidelberg New York, 1973.
 \end{thebibliography}






\end{document}
